\section{Executive Summary}

\subsection{What the project does}

LM-IDNet learns a device's ordinary network traffic from categorical counts.
It converts packet captures into ten-minute windows of TCP, ordinary UDP, SSDP,
and ARP counts, fits a Dirichlet--multinomial model on an earlier period, and
calibrates an anomaly threshold on a separate benign period. A later window
receives a score, an anomaly decision, and an explanation based on its traffic
composition, total packet volume, and ARP activity. Attack labels are used for
evaluation, not to fit the normal model or decide whether to learn from a runtime window.
This is unsupervised, normal-only learning: it still requires a trusted benign
reference. No attack-signature database or deep-learning model is involved.
Protection consists of detection and alert evidence; automatic blocking or
remediation is outside the implemented pipeline.

Batch commands produce research reports over configured captures. A streaming
command receives one completed count window at a time and immediately emits
its event. Both call the same one-window implementation. Static operation,
threshold adaptation, and periodic refitting remain available under explicit
conservative and historical legacy policies.

\begin{center}
\fbox{\begin{minipage}{0.94\textwidth}
\centering
capture bytes $\longrightarrow$ validated device windows
$\longrightarrow$ fitted normal model $\longrightarrow$ benign calibration\\[0.35em]
$\longrightarrow$ shared static/adaptive detector
$\longrightarrow$ events, accepted state, and offline evaluation
\end{minipage}}
\end{center}

\subsection{What is complete and what still needs evidence}

The agreed minimum supplied-window software is implemented, including input
compatibility, conservative learning gates, atomic state activation, restart,
rollback, bounded streaming, offline packaging, service operation, and a local
one-CPU/512-MiB Docker simulation. This is a usable software deliverable.
It does not follow that every research question has a positive answer or that
the final experimental protocol is complete.

\begin{longtable}{>{\raggedright\arraybackslash}p{0.24\textwidth}
>{\raggedright\arraybackslash}p{0.23\textwidth}
>{\raggedright\arraybackslash}p{0.43\textwidth}}
\caption{Current implementation and evidence status.}\label{tab:executive-status}\\
\toprule
\textbf{Area} & \textbf{Status} & \textbf{Evidence and limit} \\
\midrule\endfirsthead
\toprule\textbf{Area} & \textbf{Status} & \textbf{Evidence and limit} \\
\midrule\endhead
Reference detector & \working & Device windows, fitting, calibration,
explainable scoring, and labelled development results. \\
Native LM backend & \working & Numerical tests, SciPy comparisons, matrix
optimisation, and retained production/paper-style benchmarks. \\
Legacy adaptation & \working & Six explicit configurations, retained positive
and negative development results, and a historical reproduction record. \\
Conservative adaptation & \working & Admission and movement guards, controlled
update tests, and retrospective cross-device comparisons; learning benefits
remain device-dependent. \\
Persistence and rollback & \working & Versioned complete pairs, a writer lock,
atomic activation, bounded history, recovery, rollback, and failure tests. \\
Edge simulation & \working & Offline Linux/amd64 image, bounded service,
restart/stop checks, and 2,000/20,000-input-window resource trials. \\
Temporal research & \wip & Separate branch with actual-timestamp validation,
five rolling folds, and capture/episode metrics. \\
Detection and scoring & \working & DM composition with jointly calibrated
packet/ARP counts and fixed update references; broader features rejected or
excluded where common input evidence is unavailable. \\
Independent final evaluation & \notdone & Later partitions have retrospective
audits; ordinary scoring/evaluation commands select development data. \\
Empirical forecasting & \notdone & The CLI is synthetic scaffolding; a separate
paper-style experiment has a narrower forecast comparison. \\
Physical gateway validation & Outside agreed simulation & No physical/ARM
performance, packet collector, or durable event-delivery claim. \\
Thesis workspace & Working skeleton & Modular English LaTeX sources and a
tested Docker build; research prose and final evidence remain to be written. \\
\bottomrule
\end{longtable}

\subsection{The main findings}

On the Chromecast and Samsung camera development sets, static precision is
88.24\% and 92.31\%, with recall of 46.88\% and 68.57\%, respectively.
These results support high precision on those traces, with incomplete attack
coverage; they do not establish the same precision across devices.
Historical periodic refitting improves the
Chromecast development F1 from 61.22\% to 80.70\% at the selected 288/240 setting,
but the same setting does not improve Samsung over its static baseline.
The historical Samsung learning buffer also admits attacks and permits large
concentration changes. Those findings motivate the conservative learning policy
and the fixed detection references used during updates.

Across 2,038 labelled Chromecast and Samsung windows, retained scoring and
reference checks preserve static and conservative-adaptive F1 at 74.76\%.
Under legacy adaptation, pooled threshold F1 improves from 56.50\% to 74.76\%,
and periodic-refit F1 from 57.30\% to 75.36\%. Static score ranking improves on
Chromecast, while individual device/mode tradeoffs remain. These retrospective results are
distinct from evidence of useful model refitting. The long synthetic edge replay
promoted no candidate: one attempt
lacked enough accepted windows and 110 actual fits did not converge. Successful
promotion is checked separately with controlled candidates. These two kinds of
evidence answer different questions and must remain separate in the thesis.

The native LM backend preserves the intended row-wise likelihood and agrees
closely with SciPy. Its speed advantage depends on the workload; fitting and
single-window scoring do not have the same result. The resource simulation
observed about 69 MiB peak sampled process memory and a worst complete window
processing time of 5.079 seconds, including refitting, under a one-CPU quota.

The labelled 10/20/30/60-minute study finds the strongest pooled static
development F1 at thirty minutes, with 45/54 detected attacks versus 38/54 at
ten minutes. Alerts arrive by attack end for twelve versus 32 attacks,
respectively. Longer-window threshold adaptation also increases Samsung false
alerts. The reference stays at ten minutes for response time and its current
validated starting pairs; thirty minutes remains a useful static candidate.

\subsection{How to use this report}

Read the scientific and implementation sections to understand the delivered
system, the result sections to identify supported claims, and the appendices to
locate code, records, and reproduction commands. Historical branch evidence is
included because it explains decisions and failed alternatives; the retained
scoring contract and its comparisons are identified separately. The immediate
next work is a frozen, reproducible final evaluation and thesis writing; broader
feature research is optional unless the thesis scope explicitly requires it.
